% Rúbrica: Procesamiento de datos. Fuente: notebooks/00_exploracion.ipynb, bloques 1 y 2 y §3.5;
% backend/analisis/limpieza.py y construir_features de backend/modelado/entrenar_baseline.py.
\section{Procesamiento de datos}\label{sec:procesamiento}

El procesamiento tiene dos destinos. El modelo de riesgo se entrena con data-v1 tal como se
publicó, porque solo usa variables del juego y no lee texto. La limpieza produce otro conjunto,
el limpio, que usan el análisis de texto (sección~\ref{sec:exploracion-texto}) y la medición de
los motivos (sección~\ref{sec:interpretabilidad-motivos}).

\subsection{Validación y limpieza}\label{sec:procesamiento-validacion}\label{sec:procesamiento-limpieza}

\paragraph{Validación.} El notebook 00 corre primero los chequeos de la
sección~\ref{sec:datos-calidad}, y los de rango que nunca deben fallar son
aserciones\footnote{Comprobaciones dentro del código que detienen la ejecución si no se cumplen.}.
Además, \cifraTextosEntreJuegos{} textos idénticos de \cifraPalabrasDeUnaCopia{} palabras o más
aparecen en más de un juego, \cifraCopiasEnFoldsDistintos{} de ellos en folds distintos: un modelo
de texto podría verlos en entrenamiento y en prueba, y por eso la limpieza quita los duplicados.

\paragraph{Limpieza: marcar, no borrar.} Cada regla es una función de \texttt{limpieza.py}. Solo
dos quitan filas: los textos idénticos, de los que se
queda la reseña más antigua (\cifraDuplicadosQuitados{} filas), y las reseñas vacías
(\cifraVaciasQuitadas{}); quedan \cifraResenasLimpias{} reseñas. Las demás agregan una columna que
marca la reseña y la dejan en su lugar: las cortas, que son el \cifraCortasPorResena{} y cuya
salida cambiaría qué reseñas se estudian; las escritas en otro idioma, y las plantillas. La
normalización pasa el texto a minúsculas, quita el BBCode\footnote{El formato que Steam permite en
las reseñas, con etiquetas como \texttt{[h1]} o \texttt{[spoiler]}.} y las direcciones web, y abre
las contracciones para que la negación quede como palabra suelta (detalle en el notebook 00,
bloque 2). La limpieza no se aplica al modelo de riesgo. Con la misma partición congelada, sin los
duplicados su PR-AUC pasaría de \cifraPRAUCModelo{} a \cifraPRAUCSinDuplicados{}, y sin las vacías,
a \cifraPRAUCSinVacias{}; sin las cortas subiría a \cifraPRAUCSinCortas{}, pero sobre otra
población y con más prevalencia, así que no se compara. El conjunto limpio se guarda con un sha256
de su contenido.

\paragraph{Idioma.}\label{sec:procesamiento-idioma} La ingesta pidió reseñas en inglés. Según
lingua\footnote{Una biblioteca de detección de idioma que solo decide cuando tiene confianza.}, el
\cifraIdiomaInglesPorResena{} de las reseñas limpias está en inglés o no se pudo determinar, y
\cifraOtroIdioma{} están en otro idioma, sobre todo ruso, español y portugués. Como el filtro lo
aplicó Steam, esa cifra mide qué tan bien etiquetaron sus reseñas los autores. Las reseñas en otro
idioma se marcan y se quedan; al modelo no le afectan, pero a los motivos sí, porque sus palabras
clave están en inglés.

\subsection{Variables, imputación, valores extremos y aumentación}\label{sec:procesamiento-precio}

\paragraph{Ingeniería de variables.} El modelo usa cinco variables del juego. El precio entra como
el logaritmo de uno más el precio, para que los juegos caros no dominen la escala; la nota de
Metacritic entra junto con una bandera que dice si el juego la tiene, y las cinco se estandarizan
dentro del modelo (sección~\ref{sec:modelacion}).

\paragraph{Imputación.} La nota que falta se llena con la mediana de data-v1, pero siempre junto a
la bandera, así que el hueco no se rellena en silencio. Los \cifraPagoSinPrecioEntrenamiento{}
juegos de pago que llegaron sin precio entran con precio 0: el modelo los lee como si costaran 0,
aunque no los marca como gratis, y como el coeficiente del precio es positivo, eso tiende a bajar
su riesgo estimado. La ficha de esos juegos lo avisa (sección~\ref{sec:interpretabilidad}). El 0
de los perfiles privados no se imputa como biblioteca vacía: es una bandera de privacidad, y además
no entra al modelo.

\paragraph{Valores extremos.} No se recorta ninguno. Los rangos imposibles no tienen casos
(sección~\ref{sec:datos-calidad}), el precio ya entra en logaritmo, y de los minutos al reseñar
solo importa si quedan antes o después de los 120, así que una reseña escrita tras miles de horas
no pesa más que otra.

\paragraph{Aumentación de datos.} No aplica. El desbalance de la clase se atiende con pesos de clase
en el modelo y con una métrica que no lo esconde (sección~\ref{sec:objetivo-metrica}); además, lo
que varía es el juego, y crear reseñas sintéticas no agregaría juegos.

\subsection{Variables que no se conocen al comprar}\label{sec:procesamiento-variables}

Una variable solo puede entrar al modelo si se conoce antes de comprar (tabla~\ref{tab:variables}).
Las del juego vienen de la tienda y entran. Las del autor de la reseña, no:
\texttt{num\_reviews} cambia con el tiempo y se lee al descargar, y \texttt{steam\_purchase},
\texttt{received\_for\_free} y \texttt{written\_during\_early\_access} solo existen porque la
reseña ya se escribió, así que ninguna está disponible al momento de la compra. La biblioteca del
autor sí podría declararse antes; se probó, y su aporte cupo en el ruido entre folds
(sección~\ref{sec:modelacion-bmas}). En el sitio, el perfil se declara en un formulario y sirve
para la afinidad: no mueve el riesgo. El texto de la reseña es otro caso: tampoco existe al
comprar, pero además sería una fuga\footnote{Hay fuga cuando las variables de un modelo ya
contienen la respuesta que debe predecir. Sus métricas salen mejores de lo que serán en uso.},
porque la señal $Y$ sale de esas mismas reseñas, de su voto y de sus minutos, y el texto acompaña a
ese voto.

\begin{table}[htbp]
  \centering\footnotesize
  % Generado por documento/generar_figuras.py. No se edita a mano.
\begin{tabularx}{\textwidth}{>{\raggedright\arraybackslash}X >{\raggedright\arraybackslash}p{4.2cm} >{\raggedright\arraybackslash}p{4.6cm}}
\toprule
Variable & Cuándo se conoce & ¿Entra al modelo de riesgo? \\
\midrule
\texttt{es\_gratis}, precio (en logaritmo) y descuento & en la tienda, antes de comprar & sí \\
Si tiene nota de Metacritic, y la nota & en la tienda, antes de comprar & sí; la nota que falta se rellena con la mediana, junto a la bandera \\
\texttt{num\_games\_owned}: juegos del autor & en la reseña; en el sitio se declara «compras al año» & no: su aporte cabe en el ruido entre folds, y 0 es un perfil privado \\
\texttt{num\_reviews}: reseñas del autor & al descargar; cambia con el tiempo & no: no está disponible al momento de la compra \\
\texttt{steam\_purchase}, \texttt{received\_for\_free} y \texttt{written\_during\_early\_access} & en la reseña & no: no están disponibles al momento de la compra \\
\texttt{playtime\_at\_review} y \texttt{voted\_up} & en la reseña & no: definen $Y$ \\
Texto de la reseña & en la reseña & no: sería una fuga, porque $Y$ sale de esas mismas reseñas \\
\bottomrule
\end{tabularx}

  \caption{Variables candidatas: cuándo se conocen y si entran al modelo de riesgo. Fuente:
  \texttt{construir\_features} de \texttt{backend/modelado/entrenar\_baseline.py} y el
  notebook 00, §3.5.}
  \label{tab:variables}
\end{table}
